% Shared by both versions: discussion and limitations.

\iffullversion
The functionality is deliberately modest, and its liveness more so:
all activity is serialised through one output, so an honest exit's delay
is set by policy and fees, which [C] names rather than derives. Cluster mempool~\cite{clustermempool} (Bitcoin Core 31) replaced the
25-transaction chain limit with a 64-transaction cluster limit, which raises $\kappa$ for a
cooperative tree but shares the cluster with every descendant anyone
attaches. Two changes were weighed against it. A relative lock on
transitions that remove no value, so that only removals could chain, was
declined: the first conflicting spend in a block already invalidates an
honest exit, so the lock would cost throughput and remove nothing. The
tower exit \tr{EXIT}$_K^{t}$ (spec 7.2 rule 7) was adopted: a signature
over the leaf without its amount and the tree's identifier, checked from the stack, lets a
tower rebuild and rebroadcast an exit with the owner offline, which takes
the owner's liveness out of the auction without changing its price. A
hypothetical native Merkle-sum opcode (simulator only, in no tier, not
proposed here) takes \tr{EXIT}$_1$ to 252 B of script at every depth and the
exit to about $381 + 10k\vB$ for $2^k$ members (\S\ref{sec:eval}).
\else
All activity is serialised through one output, so an honest exit's delay
is set by policy and fees, which [C] names; cluster
mempool~\cite{clustermempool} raises $\kappa$ for a cooperative tree but
shares the cluster with every descendant anyone attaches, and fee
pinning is outside [C]. The price of the one-transaction exit is one
UTXO for all activity. Next: a
verified compiler, a proof that channel leaves compose, a P
model~\cite{plang}, and costing a native Merkle-sum opcode (\S\ref{sec:eval}).
\fi
